\documentclass[11pt,a4paper]{article}
\usepackage[T1]{fontenc}
\usepackage{lmodern,amsmath,amssymb,amsthm,mathtools,microtype}
\usepackage[margin=27mm]{geometry}
\usepackage[colorlinks=true,linkcolor=blue,citecolor=blue,urlcolor=blue]{hyperref}
\setlength{\emergencystretch}{2em}
\newtheorem{theorem}{Theorem}[section]
\newtheorem{lemma}[theorem]{Lemma}
\newtheorem{corollary}[theorem]{Corollary}
\newcommand{\Q}{\mathbb Q}\newcommand{\Z}{\mathbb Z}
\DeclareMathOperator{\Tr}{Tr}
\title{A cofinite obstruction to finite M\"obius coupling}
\author{}\date{17 November 2026}
\begin{document}\maketitle
\begin{abstract}
For every coefficient shift in the obstructed progressions of the
six-contact M\"obius theorem, an arithmetic progression of positive
integer inputs has no prescribed-index witness. The conclusion is
simultaneous for every fixed finite menu of integer exponents at
least four, rational multipliers, and the twelve contact symmetries.
Local nonsquare classes exclude ten branches, while a simple inert
contact excludes the remaining nonidentity branch. Open subgroups
of local elliptic groups propagate these obstructions from one input
to an entire progression. Thus adjoining finitely many exceptional
inputs cannot restore completeness. After restricting the coefficient
progression further, the same input progression works uniformly.
\end{abstract}

\section{Source and statement}

Retain the source, guards and notation of the six-contact theorem
\cite{Mobius}. In particular, let
\[
 K=\Q(\rho),\quad \rho^2=2,\quad D=(3-\rho)^2,
 \qquad E:y^2=x^3-Dx,
\]
\[
 R_0=(1+\rho,2\rho),\qquad
 P=2R_0=\left(\frac{57-28\rho}{8},
                       \frac{-660+417\rho}{32}\right).
\]
The point $P$ is non-torsion, as in that source. For a positive
integer input $n$, the prescribed root is exactly
\begin{equation}\label{eq:prescribed}
                 X=n,\qquad A=[n]P.
\end{equation}
Set $T_n=\Tr_{K/\Q}x([n]P)$ and, for an odd coefficient $B\ge3$,
\[
 a_e(n)=(n^2+B)^e+T_n^e,\qquad
 t_0=\lambda,\quad t_1=a_e(n)\lambda,
\]
\[
 F(t)=\frac{t(t+1)(2t+1)(3t+1)(4t+1)(5t+2)}{120}.
\]
Let $\mathcal M$ be the twelve rational M\"obius transformations
preserving the roots of $F$. The coupling is
\begin{equation}\label{eq:coupling}
                  a_e(n)\lambda=\phi(\lambda),
                  \qquad\phi\in\mathcal M,
\end{equation}
with a nonzero M\"obius denominator.

The detectors have
\[
 d_j=\varepsilon_j F(t_j),\quad
 w_j=\frac{d_j}{4(t_j-1)^8},\quad
 W_j=\prod_{i=1}^6(1+C_i^2/d_j),\quad
 V_j=\Tr(W_j^{-1}),\quad b_j=4w_j/V_j^2.
\]
Their tests and all point, quotient and product guards remain those
of \cite{Mobius}. The argument below uses that the tests include
$1/x(H)$ and $\overline{1/x(H-P)}$, and that
$d_j(t_j-1)V_j\ne0$. The allowed pole predicates imply
\begin{equation}\label{eq:pole}
 v_p(b_j)\text{ is not a negative odd integer}
                         \qquad(p\notin S),
\end{equation}
where $S$ is one fixed finite set of primes.

\begin{theorem}\label{thm:cofinite}
Fix a nonempty finite set $\mathcal E\subset\Z_{\ge4}$, a
nonempty finite set $\mathcal U\subset\Q^\times$, and a finite
set $S$ containing $2,3,5,7$. Let $B_0\pmod M$ be a coefficient
progression supplied by \cite[Theorem~1.1]{Mobius} for these data.
For every fixed $B\ge3$ in that progression, there is an integer
$L_B\ge1$ such that no positive input
\[
                            n\equiv1\pmod{L_B}
\]
admits a guarded rational auxiliary tuple satisfying
\eqref{eq:prescribed}, \eqref{eq:coupling}, the detector equations,
and even one of the two pole conditions \eqref{eq:pole}.
This holds simultaneously for all choices of
$e\in\mathcal E$, $\phi\in\mathcal M$, and
$\varepsilon_0,\varepsilon_1\in\mathcal U$.

The pole predicates and the retained pole index may vary from tuple
to tuple, provided their implication \eqref{eq:pole} holds outside
the same $S$. In particular these coupled prescribed-index sources
are not complete even after adjoining any finite set of inputs.
\end{theorem}

The conclusion applies to an actual defining-polynomial source if
its other equations force \eqref{eq:prescribed} at integer inputs.
An index equality only modulo an unspecified torsion translate does
not, by itself, supply this exact-root hypothesis. No assertion is
made for unrestricted root indices or for infinite menus.

\section{Local persistence and integer progressions}

\begin{lemma}\label{lem:squareclass}
Every nonzero rational nonsquare is nonsquare in some finite
completion $\Q_\ell$. Each nonzero square class in $\Q_\ell$
is open, including when $\ell=2$.
\end{lemma}
\begin{proof}
If a rational number has an odd valuation at some prime, it is
not a square in that completion. If all valuations are even,
it is a rational square times $1$ or $-1$. The nonsquare case
is $-1$ times a square, which is nonsquare in $\Q_3$.
For odd $\ell$, every element of $1+\ell\Z_\ell$ is a square;
for $\ell=2$, the same holds for $1+8\Z_2$. These assertions
follow from Hensel's lemma. Multiplying these open neighborhoods
by any fixed nonzero element proves the last assertion.
\end{proof}

\begin{lemma}\label{lem:progression}
Let $\Sigma$ be a finite set of rational primes. For each
$\ell\in\Sigma$, prescribe a neighborhood of $1$ in $\Q_\ell$
and neighborhoods of $P$ in $E(K_v)$ for every place $v\mid\ell$.
There is a positive integer $L$ such that every integer
$n\equiv1\pmod L$ belongs to all the prescribed input
neighborhoods and has $[n]P$ in all the prescribed point
neighborhoods. Good reduction is not required at these places.
\end{lemma}
\begin{proof}
The group $E(K_v)$ has arbitrarily small open subgroups at its
identity. To see this even at a bad prime or at residue
characteristic two, choose an integral Weierstrass equation.
In a sufficiently small neighborhood of the identity its formal
parameter identifies the group law with a convergent formal group
over the valuation ring\cite[Chapters IV and VII]{Silverman}.
The sets whose parameter has valuation
at least $r$, for sufficiently large $r$, are open subgroups and
form a neighborhood basis. This uses neither good reduction nor
division by the residue characteristic.

Choose such a subgroup $U_v$ with $P+U_v$ inside the desired
neighborhood. The group $E(K_v)$ is compact: its projective
model is a closed subset of the compact space
$\mathbb P^2(K_v)$. Thus the open quotient $E(K_v)/U_v$ is finite.
Choose an integer $m_v\ge1$ annihilating the image of $P$ in
this quotient. Every multiple of $[m_v]P$ lies in $U_v$.
For each $\ell$, choose $r_\ell$ such that
$1+\ell^{r_\ell}\Z_\ell$ lies in its input neighborhood.
Any common positive multiple of the finitely many $m_v$ and
$\ell^{r_\ell}$ has the required property.
\end{proof}

In particular, a finite collection of continuous local conditions
on $(n,[n]P)$ which holds at $(1,P)$ persists on an arithmetic
progression. Conditions on traces are included: at $\ell$, the
trace is the sum of the continuous maps
$\Tr_{K_v/\Q_\ell}x(A_v)$ over all $v\mid\ell$.
It is necessary to prescribe all those places, including both
factors when $\ell$ splits. At a finite affine point a rational
function whose denominator is nonzero is continuous, whether or
not that denominator is a local unit. Shrinking the neighborhood
preserves its nonvanishing and valuation.

\section{The three types of symmetry}

\begin{proof}[Proof of Theorem~\ref{thm:cofinite}]
Fix an allowed $B$. At input one,
\[
 T_1=57/4,\qquad a_e(1)=(1+B)^e+(57/4)^e.
\]
Besides the identity and $\phi_*(t)=-t/(5t+1)$, the ten
symmetries have representatives
$\phi(t)=(\alpha t+\beta_\phi)/(\gamma t+\delta)$ with
$\gamma\ne0$. Their coupling discriminants are
\begin{equation}\label{eq:discriminant}
 \Delta_{e,\phi}(n)
       =(\delta a_e(n)-\alpha)^2
                         +4\gamma\beta_\phi a_e(n).
\end{equation}
The square-gap argument of \cite[Lemma~2.2]{Mobius} shows that
each $\Delta_{e,\phi}(1)$ is a nonzero rational nonsquare.
Choose a prime $\ell_{e,\phi}$ where it remains nonsquare,
using Lemma~\ref{lem:squareclass}. Continuity and openness of
its nonzero local square class give neighborhoods of $(1,P)$
at every place above that prime on which
$\Delta_{e,\phi}$ stays nonsquare. This includes primes of
bad reduction and the prime two; none is discarded.

For the remaining involution, the coefficient progression supplies,
for each $e$, a prime $p_e\notin S$ inert in $K$, outside the
numerator and denominator supports of every element of
$\mathcal U$, such that
\begin{equation}\label{eq:contact}
                         v_{p_e}(a_e(1)-4)=1.
\end{equation}
The choices in the proof of \cite[Theorem~1.1]{Mobius} also make
$p_e\ne2,3,5,7$. Choose further neighborhoods on which
\[
                     a_e(n)-a_e(1)\in p_e^2\Z_{p_e}.
\]
This preserves \eqref{eq:contact} with $n$ in place of one.
Only finitely many places and local conditions have been imposed.
Lemma~\ref{lem:progression} provides $L_B$ satisfying all of them.

Take a positive $n\equiv1\pmod{L_B}$. The point $[n]P$ is a
nonzero double at both real embeddings. If $Q=[n]R_0$, the
duplication formula is
\[
                 x(2Q)=\left(\frac{x(Q)^2+D}{2y(Q)}\right)^2>0.
\]
Its denominator is nonzero because $Q$ is non-torsion, and $D$
is positive at both embeddings. In fact $x([n]P)>\sqrt D$
at either embedding, by the curve equation and non-torsion.
Thus $T_n>6$ and $a_e(n)>4$, including for odd $e$.
The identity coupling forces $\lambda=0$, violating $d_0\ne0$.
For each of the other ten listed branches,
$\gamma a_e(n)\ne0$, and its quadratic coupling equation has
the locally nonsquare discriminant \eqref{eq:discriminant}.
It therefore has no rational solution.

It remains to consider $\phi_*$. Since $\lambda=0$ is forbidden,
its only candidate is
\[
              \lambda=-\frac{a+1}{5a},\qquad a=a_e(n).
\]
Put
\[
 H(a)=(a^2-1)(4a-1)(3a-2)(2a-3)(a-4).
\]
The identities from \cite[Section~3]{Mobius} are
\[
 -F(\lambda)=\frac{H(a)}{375000a^6},\qquad
 F(a\lambda)=\frac{H(a)}{375000},\qquad H'(4)=11250.
\]
At $p_e$, the preserved valuation gives $a\equiv4$ and
$v_{p_e}(H(a))=1$. The displayed denominators, both
$t_j-1$, and the M\"obius denominator $5\lambda+1=-1/a$
are units there. Consequently
\[
                 v_{p_e}(d_j)=v_{p_e}(w_j)=1\qquad(j=0,1)
\]
for every allowed multiplier.

For completeness, the inert-prime implication is uniform in all
the remaining auxiliary points. At this good inert prime, $P$
reduces to an affine point rather than the identity. Therefore
$H$ and $H-P$ cannot both reduce to the identity. Conjugation
preserves the unique valuation, so at least one of
$1/x(H)$ and $\overline{1/x(H-P)}$ has nonpositive valuation.
For each defined nonzero test,
\[
 v_{p_e}(1+C_i^2/d_j)=\min\{0,2v_{p_e}(C_i)-1\}\le0;
\]
the two summands have different valuation parity. A zero test
contributes the factor one. At least one factor has valuation
at most $-1$, whence $W_j^{-1}\in p_e\mathcal O_{K_{p_e}}$
and $V_j\in p_e\Z_{p_e}$. If $V_j=0$, a guard fails.
Otherwise $v_{p_e}(b_j)=1-2v_{p_e}(V_j)$ is negative and odd.
Either pole condition is impossible. This proves the theorem.
\end{proof}

\section{Uniform shifts and finite repairs}

\begin{corollary}\label{cor:uniform}
For the same finite data, there are integers $B_*\ge3$, $M'>0$
and $L>0$, with $M'$ even, such that every pair of positive
integers with $B\ge3$ satisfying
\[
               B\equiv B_*\pmod{M'},\qquad n\equiv1\pmod L
\]
has the obstruction of Theorem~\ref{thm:cofinite}.
The progression for $B$ is a subprogression of the one in that theorem.
\end{corollary}
\begin{proof}
Choose one $B_*\ge3$ in the original progression. Apply the
preceding finite local construction at $(B_*,1,P)$, now viewing
the discriminants and contact expressions as continuous functions
of all three entries. Product neighborhoods preserve their local
conditions. Choose $M'$ divisible by the original modulus and
the finitely many prime powers needed for the $B$ neighborhoods.
Use Lemma~\ref{lem:progression} for the other two entries.
The real inequality $a_e(n)>4$ holds for every positive $n$ and
every $B\ge3$, so the same proof applies uniformly.
\end{proof}

\begin{corollary}\label{cor:repair}
Fix an obstructed $B$. Suppose a polynomial presentation of the
coupled source forces the exact root prescription
\eqref{eq:prescribed} at integer inputs. Multiplying its defining
polynomial by any nonzero polynomial in the input alone, or
adjoining any finite set of inputs by any other means, cannot
make its rational input projection contain all but finitely many
positive integers.
\end{corollary}
\begin{proof}
Theorem~\ref{thm:cofinite} excludes infinitely many positive
integers in one arithmetic progression. A nonzero univariate
polynomial has only finitely many rational roots. Multiplication
adjoins at most those inputs; any other finite insertion has the
same limitation.
\end{proof}

The input modulus in Theorem~\ref{thm:cofinite} may depend on the
fixed shift. Corollary~\ref{cor:uniform} obtains uniformity by
restricting the shifts further. These are local continuity and
finite-group arguments; no assertion of all-place local solubility
or a numerical defining-variable improvement is involved.

\begin{thebibliography}{9}
\bibitem{Mobius}
\emph{An obstruction to finite M\"obius coupling of six elliptic contacts},
project manuscript, 17 November 2026,
\nolinkurl{negative_results/papers/pole_parameters/six_contact_mobius_obstruction_20261117_v63.tex}.
\bibitem{Silverman}
J.~H.~Silverman, \emph{The Arithmetic of Elliptic Curves},
second edition, Graduate Texts in Mathematics 106,
Springer, 2009; Chapters IV and VII.
\end{thebibliography}
\end{document}
